\begin{annexitem}{annex:p12.har.kprime-bits}{thm:p12-har-kprime-bits}{Every boolean computation and both semantic directions}
\begin{proof}
Each definition is an abstraction applying the already typed data
selector choose to three nat arguments. Its type follows by those
three application rules. Actual choose computations give the complete
table (0 abbreviates the constructor O, 1 the constructor S O):
\[
\begin{array}{cc|ccc}
 a&b&\mathsf{andb}\ a\ b&\mathsf{orb}\ a\ b&\mathsf{notb}\ a\\\hline
 0&0&0&0&1\\
 0&1&0&1&1\\
 1&0&0&1&0\\
 1&1&1&1&0
\end{array}
\]
There is no omitted bit pair. Under the two Bit premises, eliminate
their disjunctions into the conjunction of the stated range and iff
propositions. In each branch, congruence for the appropriate operation
transports the two input equations to the corresponding table expression;
its constructor computation gives the output equation. This is book
equality congruence, not reflecting an input equation into conversion.

For each row, a right-hand true/false condition holds exactly for the
listed output. The required forward implication projects the appropriate
input equation(s), or eliminates an impossible output equality O=one
using constructor separation. Conversely the listed input condition(s)
select that output equation, or give the same separation contradiction
in a mismatched row. These four rows prove both andb directions, both
orb directions and each notb direction, as well as their output ranges.
The notb result uses only the a column, so either b row supplies its
two identical computations.

Finally, for Bit a, a=O contradicts a=one by equality transitivity and
the proved separation of O and S O; this proves the forward negation
direction. Conversely eliminate Bit a into a=O: its O branch returns
the given equation, while its one branch contradicts the supplied
negation and eliminates FF into that equation proposition. Exchanging
O and one proves the other displayed equivalence. All uses of the
range disjunction and FF conclude Prop; the data operations were
constructed independently of these proofs.
\end{proof}
\end{annexitem}

\begin{annexitem}{annex:p12.har.kprime-order}{thm:p12-har-kprime-order}{All comparison cases against the actual source order}
\begin{proof}
The seed l0 has type nat-to-nat. Under n:nat and q:nat-to-nat,
the inner recursor has seed O:nat and step
\(\lambda j:\mathsf{nat}.\lambda r:\mathsf{nat}.q j\),
of type nat-to-nat-to-nat. Its application to m returns nat, so
lS has type nat-to-(nat-to-nat)-to-(nat-to-nat). The outer instance
therefore has target nat-to-nat and returns the stated leb type.
The displayed three computations follow by literal expansion and the
zero/successor recursor contractions, with the recursive value ignored
in the inner step. The ltb abstraction simply applies leb to S n,m.

Induct on n with the Prop motive universally quantifying m and asserting
the leb range and both iff statements. At n=O, leb O m computes to
one for every m. The source bound O<=m is proved by nat induction
on m: reflexivity at O and \(\mathsf{leseq\_suc}\) in the step. This gives the
positive equivalence. The false-code equation one=O is impossible;
the negative order premise likewise contradicts this O<=m proof.
Both false-code implications therefore eliminate FF into their target
propositions. Its range is the one alternative.

At n=S k, split m. For m=O, the code computes to O, and
S k<=O is precisely k<O, refuted by source \(\mathsf{non\_lt0}\). Thus the
negative equivalence holds, the positive equivalence has only its
separation/ex-falso branches, and its range is the O alternative.
For m=S j, the code reduces to leb k j. The source successor
monotonicity and cancellation proved in the order library give
S k<=S j iff k<=j. Their two implications also transport the respective
negations constructively: compose a negative proof with the implication
in the opposite direction. The induction hypothesis at j therefore
gives both equivalences and the unchanged range. These exhaust m.
Finally ltb n m is leb(S n)m and source n<m means S n<=m,
so its range and both semantic directions are exactly these instances.
No classical order decision was used to construct the functions.
\end{proof}
\end{annexitem}

\begin{annexitem}{annex:p12.har.kprime-any}{thm:p12-har-kprime-any}{The full inclusive-fold induction and its negative universal}
\begin{proof}
Under N:nat,p:nat-to-nat, p O:nat and
\(\lambda k:\mathsf{nat}.\lambda r:\mathsf{nat}.
\mathsf{orb}\ r\ (p(S k))\) has type nat-to-nat-to-nat.
Thus the explicit data recursion has target nat. Its computations are
Any O p=p O and Any(S N)p=orb(Any N p)(p(S N)), by expansion
and the source constructor contraction.

Fix p and its actual all-index Bit proof. Induct on N with the Prop
motive conjoining the range, positive-existential equivalence and
negative-universal equivalence. At zero the result is p O, so the
range is the given instance. A true-code proof introduces witness O
with reflexivity O<=O. Conversely a local witness i with i<=O has
i=O by the source zero-bound lemma; congruence/Prop equality transport
sends p i=one to p O=one. This is the reverse positive implication.
For the negative universal, p O=O transports to every i<=O along
that same i=O equation. Conversely instantiate the given universal at
O with reflexivity. These are both base implications.

At S N, the old output is a Bit by the induction hypothesis and
p(S N) is a Bit by the given range premise. The boolean lemma proves
the output range. For the positive equivalence, orb is one exactly
when the old result is one or p(S N)=one. In the first branch the
induction hypothesis supplies a local witness i<=N; extend its bound
to i<=S N by the source successor constructor, then introduce the
same witness. In the second branch introduce S N with reflexivity.
Conversely a witness i<=S N splits, by the proved successor-bound
lemma, into i<=N or i=S N. The first branch gives the old output
one by the reverse induction implication and hence the new orb one.
The second transports the witness's code equality to p(S N)=one
and uses the other orb implication. All witness eliminations conclude
the new code-equation proposition.

For the negative equivalence, orb is O exactly when the old result
is O and p(S N)=O. The old result's negative induction implication
gives every zero code at i<=N. For any i<=S N the same bound split
therefore supplies either the old universal instance or the new endpoint
equation transported along i=S N. Conversely a universal for i<=S N
restricts to i<=N using successor extension; the reverse induction
implication gives the old result O. Its instance at S N with reflexivity
gives the new endpoint O. The orb false implication then proves the
new result O. This closes the full induction, including both endpoints.

The universal zero-code condition refutes a hit existential: eliminate
its witness i and two proofs into FF; the universal gives p i=O,
contradicting p i=one by constructor separation. Conversely, from
negation of the hit existential, eliminate the already proved result
range into the equation Any N p=O. Its O branch returns that equation;
its one branch supplies a hit by the positive equivalence, contradicts
the assumed negation, and eliminates FF into the equation proposition.
This proves the third equivalence without DN, and uses no Set-valued
extraction from a bounded existential.
\end{proof}
\end{annexitem}

\begin{annexitem}{annex:p12.har.kprime-proper}{thm:p12-har-kprime-proper}{The actual two bounded scans and every proper-divisor implication}
\begin{proof}
Under d,n,x:nat, times x d:nat and eqb(times x d)n:nat, and its
Bit proof is the equality-code theorem at those two actual arguments.
Thus the function supplied to Any in Coeff has the required all-index
range. Its data typing and output range follow from the bounded-fold
construction. The positive fold specification gives
\(\exists x.(x\le n)\wedge(\mathsf{eqb}(xd)n=\mathsf{one})\).
Eliminate this into the corresponding bounded-equation existential;
the equality-code forward implication gives xd=n and the same witness
is introduced. Conversely any such witness gives its eqb code one by
the reverse equality implication and then the Any true equation.
This proves both Coeff directions, without a premise that d is positive.

Both strict comparisons in Proper are Bits by the comparison theorem,
and Coeff is a Bit by the preceding paragraph. Two applications of
the boolean lemma give its range and the equivalence
\[
 \mathsf{Proper}\ d\ n=\mathsf{one}
 \leftrightarrow
 (\mathsf{one}<d)\wedge(d<n)\wedge
                     \exists x.(x\le n)\wedge(xd=n).
\]
Each forward projection is the corresponding code's proved semantic
implication; each reverse implication constructs the three true code
equations and combines them by andb's reverse true implication.
From a true code, its local bounded coefficient yields div d n by
the source ex constructor; pair it with the first two bounds to get
propdiv d n. Conversely from propdiv, the first bound implies one<=d
by source \(\mathsf{lthan\_leseq}\). Only now apply the coefficient-bound theorem
to its div proof; eliminate the returned Prop existential into the
Coeff true-equation proposition. The two other code equations follow
from the given strict bounds. Combining them proves Proper d n=one.
Thus no false general coefficient bound for d=O was imported.

The data function lambda d.Proper d n has all-index Bit by that
proof, so Any types HasProper and proves its range. Its positive
equivalence supplies a local d<=n and Proper d n=one; the latter
gives propdiv, and introducing d gives the source proper-divisor
existential. Conversely, eliminate a source proper-divisor existential
into HasProper n=one. Its witness d has d<n, hence d<=n, and
the reverse Proper implication just proved gives its code one.
Introduce this bounded hit and apply the Any implication. This proves
both HasProper directions, with the outer bound genuinely justified.

For the zero-code equivalence at either Proper or HasProper, its
range and positive equivalence suffice constructively. A zero equation
and a positive proposition would imply a one equation and then O=one,
giving its negation. Conversely eliminate its range into the zero
equation: the zero branch returns it; the one branch yields the
positive proposition, contradicts its supplied negation, and eliminates
FF into the equation. These are full negative decision proofs, not
double-negation eliminations. Every existential elimination in this
item has a Prop conclusion; the scans' data expressions were defined
before any such argument.
\end{proof}
\end{annexitem}

\begin{annexitem}{annex:p12.har.kprime-correct}{thm:p12-har-kprime-correct}{The full term composition, positive/negative adequacy and zero/one computations}
\begin{proof}
The two Any applications in Coeff/HasProper have nat bounds and
nat-to-nat predicate arguments. Proper applies only nat-to-nat-to-nat
functions to nat outputs; HasProper is nat-to-nat. Finally notb on
HasProper n and ltb one n are nat, and their andb is nat. Abstracting
over each displayed datum gives every advertised program type and in
particular KPrime:nat-to-nat. No code-range proof is an argument of
the constructed term; those proofs justify its semantics only.

HasProper n and ltb one n satisfy Bit by the preceding items, so notb
and andb give KPrime's range. Their true-code equivalences give exactly
\[
 \mathsf{KPrime}\ n=\mathsf{one}
 \leftrightarrow
 (\mathsf{one}<n)\wedge(\mathsf{HasProper}\ n=O)
 \leftrightarrow
 (\mathsf{one}<n)\wedge\neg\exists d.\mathsf{propdiv}\ d\ n.
\]
For precision, forward adequacy projects the andb true-code pair,
uses the ltb implication for the first conjunct, and the notb/HasProper
zero implications for the second. Reverse adequacy takes the two prime
conjuncts, builds the ltb true equation and HasProper zero equation,
uses notb's reverse true implication, and then andb's reverse implication.
The last formula is the actual source definition prime n.

If KPrime n=O, any prime proof gives KPrime n=one by the proved
reverse implication; equality transitivity and constructor separation
give FF. This constructs neg prime n. Conversely, given neg prime n,
eliminate the proved KPrime range into KPrime n=O: return the zero
equation in its zero branch; in its one branch use forward adequacy,
contradict the negative prime proof, and eliminate FF into the zero
equation. This supplies both negative directions without DN.
Eliminating the same range into prime n or neg prime n, using its
one or zero adequacy branch respectively, supplies constructive decision.

For input O, ltb one O expands to leb(S one)O and contracts to O.
The outer andb is choose O (notb(HasProper O)) O; choose's O
constructor contraction returns O without needing the second argument
to be evaluated. For input one, ltb one one expands to
leb(S one)(S O), contracts to leb one O, and then to O. The
same outer andb contraction returns O. Thus both claimed literal
KPrime reductions are supplied, not inferred from their book equalities.
Their source-prime refutations also follow directly: one<O contradicts
\(\mathsf{non\_lt0}\), and one<one contradicts strict irreflexivity. These are the
exact boundary cases excluded by the program's computed initial gate.

For the worked negative input 4, the literal coefficient test at d=2
is eqb(x-times-2)4. At x=O and x=one the products are respectively
0 and 2, so those eqb codes contract to O. At x=2 the product
contracts to 4 and eqb to one. The inner Any has therefore reached
one; its x=3 and x=4 steps are orb one with the remaining predicate
argument and contract to one without needing that argument evaluated.
Thus Coeff 2 4 contracts to one. The strict tests 1<2 and 2<4 also
compute one, so Proper 2 4 contracts to one. The outer Any has zero
codes at d=O,one (the first strict gate fails), then this hit at d=2;
its d=3,4 steps retain one by the same short-circuit orb computation.
HasProper 4 contracts to one, notb to O, and the final KPrime andb
to O. Separately, coefficient 2 witnesses div 2 4 by its product
computation; the bounds 1<2 and 2<4 have reflexivity/successor-constructor
proofs, so propdiv 2 4 and its existential are explicitly inhabited.
Any prime 4 proof contradicts that existential. This supplies the
negative example's proof and reduction without a kernel execution.
\end{proof}
\end{annexitem}

\begin{annexitem}{annex:p12.har.kprime-divisor-characterization}{thm:p12-har-kprime-divisor-characterization}{Zero divisors, positive coefficients and the ordinary characterization}
\begin{proof}
Take prime n and project one<n and its negative proper-divisor
existential. Fix d and a proof div d n. All coefficient eliminations
in this argument have the target proposition d=one or d=n.
If d=O, any local coefficient x has xO=n; timesaltzero gives n=O.
Transport one<n along that equation to one<O, contradicted by the
source \(\mathsf{non\_lt0}\) lemma. Ex falso proves the disjunction.
If d=one, its first alternative is reflexivity, independently of the
divisibility proof. In the remaining data case d=S(S k), lifting
O<=k twice gives one<d, and in particular O<d.

Eliminate div d n locally with coefficient x and equation xd=n.
If x=O, computation again gives n=O and the same contradiction.
If x=S y, the source times computation is xd=yd+d=d+yd.
The additive order characterization with witness yd gives d<=xd,
and transport along xd=n gives d<=n. Trichotomy of d,n now has
three branches. Its equality branch supplies d=n, the second alternative.
Its d<n branch, together with one<d and the original div proof,
constructs propdiv d n and then the forbidden existential witness d.
Its n<d branch contradicts d<=n by strict/weak composition and
irreflexivity. Both impossible branches eliminate into the equality
disjunction. These data cases cover every d and coefficient, proving
the forward universally quantified characterization.

Conversely take the two displayed characterization conjuncts.
To prove neg exists d.propdiv d n, eliminate any such existential
into FF, with local d, one<d, d<n and div d n. Its universal
divisor proof gives d=one or d=n. In the first branch transport
one<d to one<one; in the second transport d<n to n<n.
Each contradicts the source strict irreflexivity. Pair this negation
with the given one<n. This is exactly prime n. No capitalization
convention, boundedness of zero divisors or classical axiom supplies
an unproved direction of the comparison.
\end{proof}
\end{annexitem}

\begin{annexitem}{annex:p12.har.kprime-small-factor}{thm:p12-har-kprime-small-factor}{Weak product monotonicity, all factor branches and interval exclusion}
\begin{proof}
First prove weak multiplication at every c, including O. From a<=b
the source additive characterization gives a local k with a+k=b.
Distr gives bc=(a+k)c=ac+kc. The additive characterization with
witness kc proves ac<=ac+kc; transport along the equation gives
ac<=bc. The local k was eliminated into this inequality proposition.
Thus no positivity premise is needed for this weak result.

For the no-multiple certificate, suppose qd<n<(S q)d and div d n.
Eliminate the coefficient existential into FF, retaining x and xd=n.
Trichotomy of x,q gives x=q, x<q, or q<x. The first two imply
x<=q, by transported reflexivity or source \(\mathsf{lthan\_leseq}\).
Weak multiplication gives n=xd<=qd, contradicting qd<n by a
strict/weak cycle. In the last branch S q<=x; weak multiplication
gives (S q)d<=xd=n, contradicting n<(S q)d by the opposite
strict/weak cycle. All branches give FF. This proves neg div d n;
the proof did not divide by d or cancel a zero multiplier.

For the small-factor lemma, take one<n, n<(S b)(S b) and a
proper-divisor existential. Eliminate it into the claimed small
proper-divisor existential, retaining d and one<d,d<n,div d n.
Eliminate its coefficient into that same Prop target, retaining x and
xd=n. We first construct the necessary cofactor properties.
A hypothetical x=O gives n=O by computation/equality transitivity,
contradicting one<n via \(\mathsf{non\_lt0}\) after Prop transport. A hypothetical
x=one gives n=d by the unit multiplication law, contradicting d<n
by strict irreflexivity after transport. Source n0n1lt applied to
these two negative equations gives one<x. In particular O<x.
Positive-factor strict preservation multiplies one<d by x, giving
one-times-x<d-times-x. The unit law, commutativity and xd=n
transport this to x<n. Also d is a coefficient witnessing div x n,
since dx=xd=n. Pair these proofs: propdiv x n is established, not
assumed from the word cofactor.

Trichotomy of d,b gives either d<=b or b<d (equality and d<b
give the first alternative). In the first case introduce e=d with
its original proper-divisor proof. Otherwise split x,b similarly.
If x<=b, introduce e=x with the just-constructed proper-divisor
proof. In the remaining case b<d and b<x, the definitions give
S b<=d and S b<=x. Weak multiplication twice, with commutativity
between the products, gives the actual chain
\[
 (S b)(S b)\le x(S b)=(S b)x\le dx=xd=n.
\]
This contradicts n<(S b)(S b) by a strict/weak cycle.
Ex falso proves the target Prop existential. These exhaust every
factor comparison branch. Both witnesses were introduced inside a
Prop existential; no Set-valued factor extraction or factorization
procedure has been invoked.
\end{proof}
\end{annexitem}

\begin{annexitem}{annex:p12.har.kprime-61}{thm:p12-har-kprime-61}{The complete small-factor certificate and the closed inclusive-loop reduction}
\begin{proof}
In this item use \(\overline{k}:=S^kO\), not a decimal primitive
constant. The literal plus/times/eqb reductions proved in the autarkic
annex give every concrete product displayed below. The source ordering
on literal naturals is certified by reflexivity followed by the required
number of \(\mathsf{leseq\_suc}\) constructor steps: one<61 is the bound 2<=61,
and 61<64 is 62<=64. Also 8 times 8 contracts to numeral 64.
Thus the exact hypotheses of the small-factor lemma hold with n=61,b=7.

The full no-multiple certificate table is
\[
\begin{array}{c|c|c|c}
 d&q&qd& (S q)d\\\hline
 2&30&60&62\\
 3&20&60&63\\
 4&15&60&64\\
 5&12&60&65\\
 6&10&60&66\\
 7&8&56&63
\end{array}
\]
For each row qd<61<(S q)d. On the left, 60<61 is reflexivity
at 61, while 56<61 is four successor extensions from 57<=57.
On the right, 61<62 is reflexivity at 62, and the bounds 62<=63,
62<=64,62<=65,62<=66 use respectively one,two,three,four successor
extensions. Transport along the literal product computations gives
exactly the two source inequalities in each row. The no-multiple
lemma therefore refutes divisibility of 61 by each of 2,3,4,5,6,7,
not just the prime entries of this list.

For any \(e\) with \(\mathsf{one}<e\) and \(e\leq7\), the full bounded enumeration is
\(e=2\) or \(e=3\) or \(e=4\) or \(e=5\) or \(e=6\) or \(e=7\). To justify it within
the source nat regime, perform eight nested data constructor cases.
At O, one<e contradicts \(\mathsf{non\_lt0}\); at S O it contradicts strict
irreflexivity at one. The next six cases are the displayed literals,
with their reflexive equality alternatives. The remaining shape
\(S^8 k\), together with \(e\leq7\), cancels seven successor pairs to give
\(S k\leq O\), precisely \(k<O\), contradicted by \(\mathsf{non\_lt0}\ k\). This supplies
every branch, rather than assuming a numeral canonicity theorem for
arbitrary open \(e\). The case eliminations all conclude the stated Prop
disjunction.

Suppose a proper-divisor existential for 61. The small-factor lemma
gives a local e<=7 with propdiv e 61. Its one<e conjunct and the
bounded enumeration give one of the six literals. Transport its div
proof along that equality and apply the corresponding table refutation.
Each branch yields FF. Hence the proper-divisor existential is false;
pair this negation with the already supplied one<61 to prove prime 61.
This establishes the source Prime(61) assertion at the explicitly
justified historical predicate, without DN.

Finally give an actual finite reduction argument for the closed program,
not an equality-reflection shortcut from this prime proof. The leb
constructor clauses reduce on literal inputs to one exactly for the
ordinary weak numeral comparison: repeatedly contract its S/S clause,
then use its zero-first or successor/zero clause. Thus ltb has the
literal strict-comparison behavior. The earlier literal times reduction
computes x times d as its literal product, and the earlier literal
eqb reduction gives O whenever two literal numbers differ.

In the outer scan at input 61, d=O,one,61 have a false properness
comparison, so their Proper code contracts to O by the short-circuit
andb/choose clause. For each literal d from 2 through 60, the two
strict comparisons are one. To justify the inner comparisons without
assuming soundness/consistency of the entire type theory or reflecting
the proved prime proposition into computation, give the following
complete meta-level literal-integer interval table:
\[
\begin{array}{c|c|c|c}
 \text{literal d range}&q&\max(qd)&\min((q+1)d)\\\hline
 2&30&60&62\\
 3&20&60&63\\
 4&15&60&64\\
 5&12&60&65\\
 6&10&60&66\\
 7&8&56&63\\
 8&7&56&64\\
 9\text{--}10&6&60&63\\
 11\text{--}12&5&60&66\\
 13\text{--}15&4&60&65\\
 16\text{--}20&3&60&64\\
 21\text{--}30&2&60&63\\
 31\text{--}60&1&60&62
\end{array}
\]
In each grouped row the maximum is q times its largest d, and the
minimum is (q+1) times its smallest d. The displayed integer products
give qd<61<(q+1)d for every literal d in that range.
For each literal coefficient x, either x<=q, so q=x+k and
qd=xd+kd>=xd, or x=q+1+k and xd=(q+1)d+kd>61.
These are ordinary nonnegative-integer additions/products in the
literal calculation, independently of source book equality. Thus in
the first case xd<=60, and in the second xd>=62; neither is 61.
The rows cover every literal d from 2 through 60 with no gap.
The constructed literal times reduction therefore produces a numeral
different from numeral 61, and the explicit literal eqb recursion
contracts to O. This supplies each inner comparison path without
an object-theory consistency premise or equality reflection.
The inner inclusive Any fold has seed O and 61 successor
steps orb O O, each contracting to O. Thus Coeff d 61 contracts
to O, and so does Proper d 61. These descriptions cover all 62
outer candidates and all 62 inner candidates where they are needed.

Consequently the outer Any also has zero seed and 61 zero-preserving
steps, yielding HasProper 61 to O. Its notb contracts to one.
The initial ltb one 61 contracts to one by the literal comparator,
and the final andb one one contracts to one. This concatenates a
finite beta/iota/delta reduction path for the actual displayed program.
It does not promote a propositional equality to conversion: the
path is justified by the closed constructor calculations for every
listed loop input. No executed evaluator or Coq certificate is reported.
The much larger printed Pocklington/Fermat example is not evaluated
or certified by this argument.
\end{proof}
\end{annexitem}
